% Options for packages loaded elsewhere
\PassOptionsToPackage{unicode}{hyperref}
\PassOptionsToPackage{hyphens}{url}
\documentclass[
]{article}
\usepackage{xcolor}
\usepackage[margin=1in]{geometry}
\usepackage{amsmath,amssymb}
\setcounter{secnumdepth}{-\maxdimen} % remove section numbering
\usepackage{iftex}
\ifPDFTeX
  \usepackage[T1]{fontenc}
  \usepackage[utf8]{inputenc}
  \usepackage{textcomp} % provide euro and other symbols
\else % if luatex or xetex
  \usepackage{unicode-math} % this also loads fontspec
  \defaultfontfeatures{Scale=MatchLowercase}
  \defaultfontfeatures[\rmfamily]{Ligatures=TeX,Scale=1}
\fi
\usepackage{lmodern}
\ifPDFTeX\else
  % xetex/luatex font selection
\fi
% Use upquote if available, for straight quotes in verbatim environments
\IfFileExists{upquote.sty}{\usepackage{upquote}}{}
\IfFileExists{microtype.sty}{% use microtype if available
  \usepackage[]{microtype}
  \UseMicrotypeSet[protrusion]{basicmath} % disable protrusion for tt fonts
}{}
\makeatletter
\@ifundefined{KOMAClassName}{% if non-KOMA class
  \IfFileExists{parskip.sty}{%
    \usepackage{parskip}
  }{% else
    \setlength{\parindent}{0pt}
    \setlength{\parskip}{6pt plus 2pt minus 1pt}}
}{% if KOMA class
  \KOMAoptions{parskip=half}}
\makeatother
\usepackage{color}
\usepackage{fancyvrb}
\newcommand{\VerbBar}{|}
\newcommand{\VERB}{\Verb[commandchars=\\\{\}]}
\DefineVerbatimEnvironment{Highlighting}{Verbatim}{commandchars=\\\{\}}
% Add ',fontsize=\small' for more characters per line
\usepackage{framed}
\definecolor{shadecolor}{RGB}{248,248,248}
\newenvironment{Shaded}{\begin{snugshade}}{\end{snugshade}}
\newcommand{\AlertTok}[1]{\textcolor[rgb]{0.94,0.16,0.16}{#1}}
\newcommand{\AnnotationTok}[1]{\textcolor[rgb]{0.56,0.35,0.01}{\textbf{\textit{#1}}}}
\newcommand{\AttributeTok}[1]{\textcolor[rgb]{0.13,0.29,0.53}{#1}}
\newcommand{\BaseNTok}[1]{\textcolor[rgb]{0.00,0.00,0.81}{#1}}
\newcommand{\BuiltInTok}[1]{#1}
\newcommand{\CharTok}[1]{\textcolor[rgb]{0.31,0.60,0.02}{#1}}
\newcommand{\CommentTok}[1]{\textcolor[rgb]{0.56,0.35,0.01}{\textit{#1}}}
\newcommand{\CommentVarTok}[1]{\textcolor[rgb]{0.56,0.35,0.01}{\textbf{\textit{#1}}}}
\newcommand{\ConstantTok}[1]{\textcolor[rgb]{0.56,0.35,0.01}{#1}}
\newcommand{\ControlFlowTok}[1]{\textcolor[rgb]{0.13,0.29,0.53}{\textbf{#1}}}
\newcommand{\DataTypeTok}[1]{\textcolor[rgb]{0.13,0.29,0.53}{#1}}
\newcommand{\DecValTok}[1]{\textcolor[rgb]{0.00,0.00,0.81}{#1}}
\newcommand{\DocumentationTok}[1]{\textcolor[rgb]{0.56,0.35,0.01}{\textbf{\textit{#1}}}}
\newcommand{\ErrorTok}[1]{\textcolor[rgb]{0.64,0.00,0.00}{\textbf{#1}}}
\newcommand{\ExtensionTok}[1]{#1}
\newcommand{\FloatTok}[1]{\textcolor[rgb]{0.00,0.00,0.81}{#1}}
\newcommand{\FunctionTok}[1]{\textcolor[rgb]{0.13,0.29,0.53}{\textbf{#1}}}
\newcommand{\ImportTok}[1]{#1}
\newcommand{\InformationTok}[1]{\textcolor[rgb]{0.56,0.35,0.01}{\textbf{\textit{#1}}}}
\newcommand{\KeywordTok}[1]{\textcolor[rgb]{0.13,0.29,0.53}{\textbf{#1}}}
\newcommand{\NormalTok}[1]{#1}
\newcommand{\OperatorTok}[1]{\textcolor[rgb]{0.81,0.36,0.00}{\textbf{#1}}}
\newcommand{\OtherTok}[1]{\textcolor[rgb]{0.56,0.35,0.01}{#1}}
\newcommand{\PreprocessorTok}[1]{\textcolor[rgb]{0.56,0.35,0.01}{\textit{#1}}}
\newcommand{\RegionMarkerTok}[1]{#1}
\newcommand{\SpecialCharTok}[1]{\textcolor[rgb]{0.81,0.36,0.00}{\textbf{#1}}}
\newcommand{\SpecialStringTok}[1]{\textcolor[rgb]{0.31,0.60,0.02}{#1}}
\newcommand{\StringTok}[1]{\textcolor[rgb]{0.31,0.60,0.02}{#1}}
\newcommand{\VariableTok}[1]{\textcolor[rgb]{0.00,0.00,0.00}{#1}}
\newcommand{\VerbatimStringTok}[1]{\textcolor[rgb]{0.31,0.60,0.02}{#1}}
\newcommand{\WarningTok}[1]{\textcolor[rgb]{0.56,0.35,0.01}{\textbf{\textit{#1}}}}
\usepackage{graphicx}
\makeatletter
\newsavebox\pandoc@box
\newcommand*\pandocbounded[1]{% scales image to fit in text height/width
  \sbox\pandoc@box{#1}%
  \Gscale@div\@tempa{\textheight}{\dimexpr\ht\pandoc@box+\dp\pandoc@box\relax}%
  \Gscale@div\@tempb{\linewidth}{\wd\pandoc@box}%
  \ifdim\@tempb\p@<\@tempa\p@\let\@tempa\@tempb\fi% select the smaller of both
  \ifdim\@tempa\p@<\p@\scalebox{\@tempa}{\usebox\pandoc@box}%
  \else\usebox{\pandoc@box}%
  \fi%
}
% Set default figure placement to htbp
\def\fps@figure{htbp}
\makeatother
\setlength{\emergencystretch}{3em} % prevent overfull lines
\providecommand{\tightlist}{%
  \setlength{\itemsep}{0pt}\setlength{\parskip}{0pt}}
\usepackage{bookmark}
\IfFileExists{xurl.sty}{\usepackage{xurl}}{} % add URL line breaks if available
\urlstyle{same}
\hypersetup{
  pdftitle={Practice Assignment},
  pdfauthor={Jonathan Dale},
  hidelinks,
  pdfcreator={LaTeX via pandoc}}

\title{Practice Assignment}
\author{Jonathan Dale}
\date{2026-02-09}

\begin{document}
\maketitle

\textbf{Instructions:} Create a \textbf{\emph{pdf}} file with all
relevant code and outputs using \textbf{Rmarkdown}. Do not submit any
other file formats and do not print the entire dataset with the answers.
\emph{ONLY include relevant outputs.}

\subsubsection{Q1) Use the `carseats.csv' for the data set. Read the
description of the dataset and answer the following
questions.}\label{q1-use-the-carseats.csv-for-the-data-set.-read-the-description-of-the-dataset-and-answer-the-following-questions.}

\begin{enumerate}
\def\labelenumi{\alph{enumi})}
\tightlist
\item
  Import the ``carseats'' dataset, look at the first few rows and
  inspect the data types of the variables in the dataframe.
\item
  Change the variables ``ShelveLoc, urban, US'' into a factor variables.
\item
  create a new variable called ``profit'' which stands for ``Income -
  Advertising''
\item
  Check for missing data. If you have missing data remove the
  corresponding rows from the dataset.
\item
  How many ``Good'' shelving locations are there in the dataset?
\item
  How many stores are inside the USA? create a separate data frame
  containing all stores from USA. Name the data set as ``stores\_USA''
\item
  create another data set called ``HighUrban\_USSales'' using
  `stores\_USA' data set. Where, sales are greater than 7 thousand and
  stores are located in Urban areas.
\item
  Remove ``US'' and ``Urban'' columns from the ``HighUrban\_USSales''
  dataset.
\item
  For one the above subset, write to a new CSV file
\end{enumerate}

\begin{center}\rule{0.5\linewidth}{0.5pt}\end{center}

\begin{Shaded}
\begin{Highlighting}[]
\FunctionTok{library}\NormalTok{(readr)}
\DocumentationTok{\#\# a) Import the "carseats" dataset, look at the first few rows and inspect the data types of the variables in the dataframe}
\NormalTok{carseats\_1 }\OtherTok{\textless{}{-}} \FunctionTok{read.csv}\NormalTok{(}\StringTok{"\textasciitilde{}/STATS/STAT{-}515{-}002/homework/hw{-}1/carseats{-}1.csv"}\NormalTok{)}
\FunctionTok{str}\NormalTok{(carseats\_1)}
\end{Highlighting}
\end{Shaded}

\begin{verbatim}
## 'data.frame':    400 obs. of  11 variables:
##  $ Sales      : num  9.5 11.22 10.06 7.4 4.15 ...
##  $ CompPrice  : int  138 111 113 117 141 124 115 136 132 132 ...
##  $ Income     : int  73 48 35 100 64 113 105 81 110 113 ...
##  $ Advertising: int  11 16 10 4 3 13 0 15 0 0 ...
##  $ Population : int  276 260 269 466 340 501 45 425 108 131 ...
##  $ Price      : int  120 83 80 97 128 72 108 120 124 124 ...
##  $ ShelveLoc  : chr  "Bad" "Good" "Medium" "Medium" ...
##  $ Age        : int  42 65 59 55 38 78 71 67 76 76 ...
##  $ Education  : int  17 10 12 14 13 16 15 10 10 17 ...
##  $ Urban      : chr  "Yes" "Yes" "Yes" "Yes" ...
##  $ US         : chr  "Yes" "Yes" "Yes" "Yes" ...
\end{verbatim}

\begin{Shaded}
\begin{Highlighting}[]
\FunctionTok{head}\NormalTok{(carseats\_1)}
\end{Highlighting}
\end{Shaded}

\begin{verbatim}
##   Sales CompPrice Income Advertising Population Price ShelveLoc Age Education
## 1  9.50       138     73          11        276   120       Bad  42        17
## 2 11.22       111     48          16        260    83      Good  65        10
## 3 10.06       113     35          10        269    80    Medium  59        12
## 4  7.40       117    100           4        466    97    Medium  55        14
## 5  4.15       141     64           3        340   128       Bad  38        13
## 6 10.81       124    113          13        501    72       Bad  78        16
##   Urban  US
## 1   Yes Yes
## 2   Yes Yes
## 3   Yes Yes
## 4   Yes Yes
## 5   Yes  No
## 6    No Yes
\end{verbatim}

\begin{Shaded}
\begin{Highlighting}[]
\DocumentationTok{\#\# b)  Change the variables "ShelveLoc, urban, US" into a factor variables.}

    \CommentTok{\# View the structure first}
  \FunctionTok{str}\NormalTok{(carseats\_1)}
\end{Highlighting}
\end{Shaded}

\begin{verbatim}
## 'data.frame':    400 obs. of  11 variables:
##  $ Sales      : num  9.5 11.22 10.06 7.4 4.15 ...
##  $ CompPrice  : int  138 111 113 117 141 124 115 136 132 132 ...
##  $ Income     : int  73 48 35 100 64 113 105 81 110 113 ...
##  $ Advertising: int  11 16 10 4 3 13 0 15 0 0 ...
##  $ Population : int  276 260 269 466 340 501 45 425 108 131 ...
##  $ Price      : int  120 83 80 97 128 72 108 120 124 124 ...
##  $ ShelveLoc  : chr  "Bad" "Good" "Medium" "Medium" ...
##  $ Age        : int  42 65 59 55 38 78 71 67 76 76 ...
##  $ Education  : int  17 10 12 14 13 16 15 10 10 17 ...
##  $ Urban      : chr  "Yes" "Yes" "Yes" "Yes" ...
##  $ US         : chr  "Yes" "Yes" "Yes" "Yes" ...
\end{verbatim}

\begin{Shaded}
\begin{Highlighting}[]
    \CommentTok{\# Adjust the variables, convert into factor variable type}
\NormalTok{  carseats\_1}\SpecialCharTok{$}\NormalTok{ShelveLoc }\OtherTok{=} \FunctionTok{factor}\NormalTok{(carseats\_1}\SpecialCharTok{$}\NormalTok{ShelveLoc)}
\NormalTok{  carseats\_1}\SpecialCharTok{$}\NormalTok{Urban }\OtherTok{=} \FunctionTok{factor}\NormalTok{(carseats\_1}\SpecialCharTok{$}\NormalTok{Urban)}
\NormalTok{  carseats\_1}\SpecialCharTok{$}\NormalTok{US }\OtherTok{=} \FunctionTok{factor}\NormalTok{(carseats\_1}\SpecialCharTok{$}\NormalTok{US)}
    \CommentTok{\# Verify that the variables were converted into factor variables}
  \FunctionTok{str}\NormalTok{(carseats\_1)}
\end{Highlighting}
\end{Shaded}

\begin{verbatim}
## 'data.frame':    400 obs. of  11 variables:
##  $ Sales      : num  9.5 11.22 10.06 7.4 4.15 ...
##  $ CompPrice  : int  138 111 113 117 141 124 115 136 132 132 ...
##  $ Income     : int  73 48 35 100 64 113 105 81 110 113 ...
##  $ Advertising: int  11 16 10 4 3 13 0 15 0 0 ...
##  $ Population : int  276 260 269 466 340 501 45 425 108 131 ...
##  $ Price      : int  120 83 80 97 128 72 108 120 124 124 ...
##  $ ShelveLoc  : Factor w/ 4 levels "","Bad","Good",..: 2 3 4 4 2 2 4 3 4 4 ...
##  $ Age        : int  42 65 59 55 38 78 71 67 76 76 ...
##  $ Education  : int  17 10 12 14 13 16 15 10 10 17 ...
##  $ Urban      : Factor w/ 3 levels "","No","Yes": 3 3 3 3 3 2 3 3 2 2 ...
##  $ US         : Factor w/ 2 levels "No","Yes": 2 2 2 2 1 2 1 2 1 2 ...
\end{verbatim}

\begin{Shaded}
\begin{Highlighting}[]
\CommentTok{\# c)  create a new variable called "profit" which stands for "Income {-} Advertising"}
  
    \CommentTok{\# Create the new variable}
\NormalTok{  profit }\OtherTok{=}\NormalTok{ carseats\_1}\SpecialCharTok{$}\NormalTok{Income }\SpecialCharTok{{-}}\NormalTok{ carseats\_1}\SpecialCharTok{$}\NormalTok{Advertising}
    \CommentTok{\# Then add the new variable to the data set}
\NormalTok{  carseats\_1}\SpecialCharTok{$}\NormalTok{profit }\OtherTok{=}\NormalTok{ profit}
    \CommentTok{\# verify}
  \FunctionTok{str}\NormalTok{(carseats\_1)}
\end{Highlighting}
\end{Shaded}

\begin{verbatim}
## 'data.frame':    400 obs. of  12 variables:
##  $ Sales      : num  9.5 11.22 10.06 7.4 4.15 ...
##  $ CompPrice  : int  138 111 113 117 141 124 115 136 132 132 ...
##  $ Income     : int  73 48 35 100 64 113 105 81 110 113 ...
##  $ Advertising: int  11 16 10 4 3 13 0 15 0 0 ...
##  $ Population : int  276 260 269 466 340 501 45 425 108 131 ...
##  $ Price      : int  120 83 80 97 128 72 108 120 124 124 ...
##  $ ShelveLoc  : Factor w/ 4 levels "","Bad","Good",..: 2 3 4 4 2 2 4 3 4 4 ...
##  $ Age        : int  42 65 59 55 38 78 71 67 76 76 ...
##  $ Education  : int  17 10 12 14 13 16 15 10 10 17 ...
##  $ Urban      : Factor w/ 3 levels "","No","Yes": 3 3 3 3 3 2 3 3 2 2 ...
##  $ US         : Factor w/ 2 levels "No","Yes": 2 2 2 2 1 2 1 2 1 2 ...
##  $ profit     : int  62 32 25 96 61 100 105 66 110 113 ...
\end{verbatim}

\begin{Shaded}
\begin{Highlighting}[]
\CommentTok{\# d)  Check for missing data. If you have missing data remove the corresponding rows from the dataset.}
  
    \CommentTok{\# Do we have missing data?}
  \FunctionTok{table}\NormalTok{(}\FunctionTok{is.na}\NormalTok{(carseats\_1)) }\DocumentationTok{\#\# Yes we have missing data!}
\end{Highlighting}
\end{Shaded}

\begin{verbatim}
## 
## FALSE  TRUE 
##  4797     3
\end{verbatim}

\begin{Shaded}
\begin{Highlighting}[]
    \CommentTok{\# Create new dataset by removing the missing values from the original.}
\NormalTok{  carseats\_clean }\OtherTok{=} \FunctionTok{na.omit}\NormalTok{(carseats\_1)}


\CommentTok{\# e)  How many "Good" shelving locations are there in the dataset?}
    
    \CommentTok{\# Take a subset with only the ShelveLoc = "Good"}
\NormalTok{  Good\_location }\OtherTok{=} \FunctionTok{subset}\NormalTok{(carseats\_1, ShelveLoc }\SpecialCharTok{==} \StringTok{"Good"}\NormalTok{)}
    \CommentTok{\# Use nrow() to count the new subset}
  \FunctionTok{nrow}\NormalTok{(Good\_location)}
\end{Highlighting}
\end{Shaded}

\begin{verbatim}
## [1] 85
\end{verbatim}

\begin{Shaded}
\begin{Highlighting}[]
    \CommentTok{\# Alternatively, just apply a filter and wrap with the nrow() function.}
  \FunctionTok{nrow}\NormalTok{(}\FunctionTok{subset}\NormalTok{(carseats\_1, ShelveLoc }\SpecialCharTok{==} \StringTok{"Good"}\NormalTok{))}
\end{Highlighting}
\end{Shaded}

\begin{verbatim}
## [1] 85
\end{verbatim}

\begin{Shaded}
\begin{Highlighting}[]
\CommentTok{\# f)  How many stores are inside the USA? create a separate data frame containing all stores from USA. Name the data set as "stores\_USA"}

    \CommentTok{\# Create new dataframe as a subset of original}
\NormalTok{  stores\_USA }\OtherTok{=} \FunctionTok{subset}\NormalTok{(carseats\_1, US }\SpecialCharTok{==} \StringTok{"Yes"}\NormalTok{)}
    \CommentTok{\# use nrow() function to count rows of the subset where US == "Yes"}
  \FunctionTok{nrow}\NormalTok{(stores\_USA)}
\end{Highlighting}
\end{Shaded}

\begin{verbatim}
## [1] 258
\end{verbatim}

\begin{Shaded}
\begin{Highlighting}[]
\CommentTok{\# g)  create another data set called "HighUrban\_USSales" using \textquotesingle{}stores\_USA\textquotesingle{} data set. Where, sales are greater than 7 thousand and stores are located in Urban areas.}
  
    \CommentTok{\# Create new data frame}
\NormalTok{  HighUrban\_USSales }\OtherTok{=} \FunctionTok{subset}\NormalTok{(stores\_USA, Sales }\SpecialCharTok{\textgreater{}} \DecValTok{7} \SpecialCharTok{\&}\NormalTok{ Urban }\SpecialCharTok{==} \StringTok{"Yes"}\NormalTok{)}
  

\CommentTok{\# h)  Remove "US" and "Urban" columns from the "HighUrban\_USSales" dataset.}
  
    \CommentTok{\# Check the original structure, remove the columns, verify  }
  \FunctionTok{str}\NormalTok{(HighUrban\_USSales)}
\end{Highlighting}
\end{Shaded}

\begin{verbatim}
## 'data.frame':    104 obs. of  12 variables:
##  $ Sales      : num  9.5 11.2 10.1 7.4 11.8 ...
##  $ CompPrice  : int  138 111 113 117 136 117 115 107 147 129 ...
##  $ Income     : int  73 48 35 100 81 94 28 117 74 76 ...
##  $ Advertising: int  11 16 10 4 15 4 11 11 13 16 ...
##  $ Population : int  276 260 269 466 425 503 29 148 251 58 ...
##  $ Price      : int  120 83 80 97 120 94 86 118 131 121 ...
##  $ ShelveLoc  : Factor w/ 4 levels "","Bad","Good",..: 2 3 4 4 3 3 3 3 3 4 ...
##  $ Age        : int  42 65 59 55 67 50 53 52 52 69 ...
##  $ Education  : int  17 10 12 14 10 13 18 18 10 12 ...
##  $ Urban      : Factor w/ 3 levels "","No","Yes": 3 3 3 3 3 3 3 3 3 3 ...
##  $ US         : Factor w/ 2 levels "No","Yes": 2 2 2 2 2 2 2 2 2 2 ...
##  $ profit     : int  62 32 25 96 66 90 17 106 61 60 ...
\end{verbatim}

\begin{Shaded}
\begin{Highlighting}[]
\NormalTok{  HighUrban\_USSales }\OtherTok{=} \FunctionTok{subset}\NormalTok{(HighUrban\_USSales, }\AttributeTok{select =} \SpecialCharTok{{-}}\FunctionTok{c}\NormalTok{(US, Urban))}
  \FunctionTok{str}\NormalTok{(HighUrban\_USSales)}
\end{Highlighting}
\end{Shaded}

\begin{verbatim}
## 'data.frame':    104 obs. of  10 variables:
##  $ Sales      : num  9.5 11.2 10.1 7.4 11.8 ...
##  $ CompPrice  : int  138 111 113 117 136 117 115 107 147 129 ...
##  $ Income     : int  73 48 35 100 81 94 28 117 74 76 ...
##  $ Advertising: int  11 16 10 4 15 4 11 11 13 16 ...
##  $ Population : int  276 260 269 466 425 503 29 148 251 58 ...
##  $ Price      : int  120 83 80 97 120 94 86 118 131 121 ...
##  $ ShelveLoc  : Factor w/ 4 levels "","Bad","Good",..: 2 3 4 4 3 3 3 3 3 4 ...
##  $ Age        : int  42 65 59 55 67 50 53 52 52 69 ...
##  $ Education  : int  17 10 12 14 10 13 18 18 10 12 ...
##  $ profit     : int  62 32 25 96 66 90 17 106 61 60 ...
\end{verbatim}

\begin{Shaded}
\begin{Highlighting}[]
\CommentTok{\# i)  For one the above subset, write to a new CSV file}
    \DocumentationTok{\#\# Write to file removing the row column}
  \FunctionTok{write.csv}\NormalTok{(HighUrban\_USSales, }\StringTok{"HighUrban\_USSales.csv"}\NormalTok{, }\AttributeTok{row.names =}\NormalTok{ F)}
\end{Highlighting}
\end{Shaded}

\subsubsection{Q2) See the following code of a function and explain what
it does. Suggest a suitable name for the function and rename.
Demonstrate how the function works when you have numerical data and
character
data.}\label{q2-see-the-following-code-of-a-function-and-explain-what-it-does.-suggest-a-suitable-name-for-the-function-and-rename.-demonstrate-how-the-function-works-when-you-have-numerical-data-and-character-data.}

The code below defines a function named `function1' that takes a single
argument x. If the vector, x, has length 1 or less, the function returns
NULL. If length greater than 1, the function returns a new vector with
all the elements except the last one. I would re-name this to
remove\_last\_element() to better describe what it dose.

\begin{Shaded}
\begin{Highlighting}[]
\NormalTok{function1 }\OtherTok{\textless{}{-}} \ControlFlowTok{function}\NormalTok{(x) \{}
  \ControlFlowTok{if}\NormalTok{ (}\FunctionTok{length}\NormalTok{(x) }\SpecialCharTok{\textless{}=} \DecValTok{1}\NormalTok{) }\FunctionTok{return}\NormalTok{(}\ConstantTok{NULL}\NormalTok{)}
\NormalTok{  x[}\SpecialCharTok{{-}}\FunctionTok{length}\NormalTok{(x)]}
\NormalTok{\}}
\end{Highlighting}
\end{Shaded}

\subsubsection{Q3) Write a function to compute the sample variance of a
numerical vector. Use the equation of the variance to write the
function.}\label{q3-write-a-function-to-compute-the-sample-variance-of-a-numerical-vector.-use-the-equation-of-the-variance-to-write-the-function.}

\begin{Shaded}
\begin{Highlighting}[]
\NormalTok{sample\_variance }\OtherTok{=} \ControlFlowTok{function}\NormalTok{(x) \{}
  \CommentTok{\# Validate the input}
  \ControlFlowTok{if}\NormalTok{ (}\SpecialCharTok{!} \FunctionTok{is.numeric}\NormalTok{(x)) \{}
    \FunctionTok{stop}\NormalTok{(}\StringTok{"Must be numeric vector"}\NormalTok{)}
\NormalTok{  \}}
  \ControlFlowTok{if}\NormalTok{ (}\FunctionTok{length}\NormalTok{(x) }\SpecialCharTok{\textless{}} \DecValTok{2}\NormalTok{) \{}
    \FunctionTok{stop}\NormalTok{(}\StringTok{"Input vector must have at least 2 elements to compute sample variance."}\NormalTok{)}
\NormalTok{  \}}
  \CommentTok{\# Input validation sources from Google AI overview.}
  
  \CommentTok{\# Calculate number of observations, mean, sum of squared differences, and variance}
\NormalTok{  n }\OtherTok{\textless{}{-}} \FunctionTok{length}\NormalTok{(x)}
\NormalTok{  xbar }\OtherTok{\textless{}{-}} \FunctionTok{mean}\NormalTok{(x, }\AttributeTok{na.rm =} \ConstantTok{TRUE}\NormalTok{)}
\NormalTok{  sum\_sq\_diff }\OtherTok{=} \FunctionTok{sum}\NormalTok{((x }\SpecialCharTok{{-}}\NormalTok{ xbar)}\SpecialCharTok{\^{}}\DecValTok{2}\NormalTok{)}
\NormalTok{  sample\_variance }\OtherTok{=}\NormalTok{ sum\_sq\_diff }\SpecialCharTok{/}\NormalTok{ (n }\SpecialCharTok{{-}}\DecValTok{1}\NormalTok{)}
  \FunctionTok{return}\NormalTok{(sample\_variance)}
\NormalTok{\}}


\NormalTok{my\_sample\_vector }\OtherTok{=} \FunctionTok{c}\NormalTok{(}\DecValTok{10}\NormalTok{, }\DecValTok{12}\NormalTok{, }\DecValTok{15}\NormalTok{, }\DecValTok{11}\NormalTok{, }\DecValTok{14}\NormalTok{, }\DecValTok{13}\NormalTok{) }\CommentTok{\# Create a sample vector}
\NormalTok{my\_sample\_variance\_function\_call }\OtherTok{=} \FunctionTok{sample\_variance}\NormalTok{(my\_sample\_vector) }\CommentTok{\# Call the function, store results in variable}
\FunctionTok{cat}\NormalTok{(}\StringTok{"Manual variance function:"}\NormalTok{, my\_sample\_variance\_function\_call, }\StringTok{"}\SpecialCharTok{\textbackslash{}n}\StringTok{"}\NormalTok{ )}
\end{Highlighting}
\end{Shaded}

\begin{verbatim}
## Manual variance function: 3.5
\end{verbatim}

\begin{Shaded}
\begin{Highlighting}[]
\DocumentationTok{\#\# Lets verify the results with the built in function...}
\NormalTok{built\_in\_sample\_variance }\OtherTok{=} \FunctionTok{var}\NormalTok{(my\_sample\_vector)}
\FunctionTok{cat}\NormalTok{(}\StringTok{"Built in sample variance var():"}\NormalTok{, built\_in\_sample\_variance, }\StringTok{"}\SpecialCharTok{\textbackslash{}n}\StringTok{"}\NormalTok{)}
\end{Highlighting}
\end{Shaded}

\begin{verbatim}
## Built in sample variance var(): 3.5
\end{verbatim}

\end{document}
